\article{关于 p 进半单群超尖点表示的若干注记}{Some Remarks on the Supercuspidal Representations of p-adic Semisimple Groups}{G. Lusztig}{1}{171--175}
\origsection{1}{一个可能的几何构造}
设 $G$ 是非阿基米德局部域（local nonarchimedean field）$K$ 上的半单群（semisimple group），$K$ 的整数环为 $\mathcal O$，有限剩余域为 $k$。定义 $\mathsf G=G(K)$、$\widetilde G=G(\widetilde K)$，其中 $\widetilde K$ 是 $K$ 的极大非分歧扩张（maximal unramified extension），其整数环为 $\widetilde{\mathcal O}$，剩余域 $\widetilde k$ 是 $k$ 的代数闭包。

设 $T\subset G$ 是在 $K$ 上定义且各向异性的（anisotropic）极大环面（maximal torus）。假定 $T$ 满足下述条件：存在包含 $T$ 且在 $\widetilde K$ 上定义的 $G$ 的 Borel 子群（Borel subgroup）$B$。（如果 $T$ 在 $\widetilde K$ 上分裂，此条件当然成立。）令 $\widehat{\mathsf T}$ 为从 $\mathsf T=T(K)$ 到 $\overline\Q_\ell^*$、通过 $\mathsf T$ 的某个有限商分解的同态集合，其中 $\ell$ 为不等于 $\operatorname{char}k$ 的素数。我们期望，对每个满足某种正则性条件（regularity condition）的 $\theta\in\widehat{\mathsf T}$，都对应一个 $\mathsf G$ 的不可约容许超尖点表示（irreducible admissible supercuspidal representation），系数域为 $\overline\Q_\ell$。Gérardin 利用 Shintani、Howe 和 Corwin 的方法，在 $T$ 于 $\widetilde K$ 上分裂且满足若干限制的情形，建立了这种对应，见 [2]。（对于 $\SL_2$，这一对应由 Gelfand、Graev、Shalika 等人建立。）

下面提出构造这一对应的另一条可能途径。这一方法拟使用 $\widetilde k$ 上某个无限维簇（infinite dimensional variety）$X$ 的 $\ell$ 进上同调（$\ell$-adic cohomology）或同调（homology）。（Drinfeld 的工作，以及 Deligne 与我在有限域情形中的工作 [1]，为利用 $\ell$ 进上同调构造 $\mathsf G$ 的表示提供了依据。在 Gérardin 的方法中，处理导子（conductor）很小的情形时，也需借助有限域上约化群的表示论，而后者本身就是建立在 $\ell$ 进上同调之上的。）

令 $U$ 为 $B$ 的幺幂根基（unipotent radical），$\widetilde U=U(\widetilde K)$。考虑由 Frobenius 元 $\phi\in\Gal(\widetilde K/K)$ 定义的 Frobenius 映射 $F:\widetilde G\to\widetilde G$，于是 $\mathsf G=\widetilde G^F$。令
\[
X=\{g\in\widetilde G\mid g^{-1}F(g)\in\widetilde U\}/(\widetilde U\cap F^{-1}\widetilde U).
\]
（$\widetilde U\cap F^{-1}\widetilde U$ 通过右乘作用于 $\widetilde G$。）于是 $\mathsf G\times\mathsf T$ 通过
\[
(g_0,t):g\longmapsto g_0gt^{-1},\qquad g_0\in\mathsf G,\ t\in\mathsf T,\ g\in\widetilde G,
\]
作用于 $X$。

我预期，可以将 $X$ 视为 $\widetilde k$ 上的无限维簇，并定义其 $\ell$ 进同调群 $H_i(X)$，使其带有 $\mathsf G\times\mathsf T$ 的作用，且满足
\[
H_i(X)=\bigoplus_{\theta\in\widehat{\mathsf T}}H_i(X)_\theta,
\]
其中 $H_i(X)_\theta$ 是 $H_i(X)$ 中 $\mathsf T$ 按 $\theta$ 作用的子空间。此外，固定 $\theta$ 后，$i$ 足够大时 $H_i(X)_\theta$ 应当为零；而若 $\theta$ 不被 $T$ 的小 Weyl 群（small Weyl group）的任何非平凡元素固定，则应当恰有一个指标 $i=i_0$（依赖于 $\theta$）使 $H_i(X)_\theta$ 非零，并且 $\mathsf G$ 在 $H_{i_0}(X)_\theta$ 上的表示应当是不可约容许超尖点表示。这样便可建立所需的 $\mathsf T$ 的特征（characters）与 $\mathsf G$ 的表示之间的对应。
\footnotetext{AMS（MOS）学科分类（1970）：主要分类 22E50；次要分类 14F20、20H30。}

\origsection{2}{中心除代数的情形}
先考虑 $\widetilde G=\SL_n(\widetilde K)$ 这一例子，令 $F:\widetilde G\to\widetilde G$ 为下式定义的同态：
\[
F(A)=\begin{pmatrix}
0&1&&0\\
&0&\ddots&\\
&&\ddots&1\\
\pi&&&0
\end{pmatrix}A^\phi
\begin{pmatrix}
0&1&&0\\
&0&\ddots&\\
&&\ddots&1\\
\pi&&&0
\end{pmatrix}^{-1}.
\tag{2.1}
\]
其中 $\pi$ 是 $\mathcal O$ 的统一化元（uniformizing element），对 $\widetilde K$ 上的矩阵 $A$，$A^\phi$ 是对 $A$ 的每个矩阵元施加 $\phi$ 所得的矩阵。不动点集（fixed point set）$\widetilde G^F$ 是 $K$ 上维数为 $n^2$ 的某个中心除代数（central division algebra）中约化范数（reduced norm）为一的元素组成的群。令 $\widetilde T$ 为 $\SL_n(\widetilde K)$ 中的对角矩阵群；它在 $F$ 下不变。令 $\widetilde U\subset\SL_n(\widetilde K)$（相应地，$\widetilde U^-$）为对角线上各元素均为 $1$ 的全部上（下）三角矩阵构成的子群。若矩阵 $A\in\SL_n(\widetilde K)$ 满足 $A^{-1}F(A)\in\widetilde U$，则存在唯一的 $B\in\widetilde U\cap F^{-1}\widetilde U$，使 $(AB)^{-1}F(AB)\in\widetilde U\cap F\widetilde U^-$。因此，
\[
\{A\in\SL_n(\widetilde K)\mid A^{-1}F(A)\in\widetilde U\}/(\widetilde U\cap F^{-1}\widetilde U)
\]
可与以下集合对应：
\[
X=\{A\in\SL_n(\widetilde K)\mid A^{-1}F(A)\in\widetilde U\cap F\widetilde U^-\}.
\]
$\widetilde G^F$ 通过左乘作用于它，$\widetilde T^F$ 通过右乘作用于它。$X$ 恰好是形如下式且行列式为一的全部 $n\times n$ 矩阵的集合：
\[
\begin{pmatrix}
a_1&a_n^{\phi^{-1}}&\cdots&a_2^{\phi^{-(n-1)}}\\
\pi a_2&a_1^{\phi^{-1}}&\ddots&\vdots\\
\vdots&\ddots&\ddots&a_n^{\phi^{-(n-1)}}\\
\pi a_n&\cdots&\pi a_2^{\phi^{-(n-2)}}&a_1^{\phi^{-(n-1)}}
\end{pmatrix},\qquad a_i\in\widetilde K.
\tag{2.2}
\]
这一矩阵的行列式条件自动蕴含 $a_i\in\widetilde{\mathcal O}$（$1\leq i\leq n$）且 $a_1\notin\pi\widetilde{\mathcal O}$。

因此，$X$ 可视为 $\widetilde k$ 上代数簇 $X_h$ 的射影极限 $\varprojlim_hX_h$。\footnote{\textbf{译者注（本次校订）：}按开头的定义，$\widetilde K$ 是未完备的极大非分歧扩张；但 $\varprojlim_h\widetilde{\mathcal O}/\pi^h\widetilde{\mathcal O}$ 是其整数环的完备化。因而这里把无限层点集直接等同于有限层的逆极限时，应在 $\widetilde K$ 的完备化中取点。第3节的同类逆极限表述也需采用这一约定。}这里 $X_h$（$h\geq1$）是所有形如（2.2）、行列式为一的 $n\times n$ 矩阵的集合，其中
\[
a_1\in(\widetilde{\mathcal O}/\pi^h\widetilde{\mathcal O})^*,\qquad
a_i\in\widetilde{\mathcal O}/\pi^{h-1}\widetilde{\mathcal O}\quad(2\leq i\leq n).
\]
（将 $\pi a_i$，$2\leq i\leq n$，视为 $\pi\widetilde{\mathcal O}/\pi^h\widetilde{\mathcal O}$ 的元素；这类矩阵的行列式可自然地看作 $\widetilde{\mathcal O}/\pi^h\widetilde{\mathcal O}$ 中的元素。）

令 $G_h$ 为行列式等于一的全部 $n\times n$ 矩阵 $(a_{ij})$ 的集合，矩阵元满足
\[
\begin{cases}
a_{ij}\in\pi\widetilde{\mathcal O}/\pi^h\widetilde{\mathcal O},&i>j,\\
a_{ij}\in\widetilde{\mathcal O}/\pi^{h-1}\widetilde{\mathcal O},&i<j,\\
a_{ii}\in(\widetilde{\mathcal O}/\pi^h\widetilde{\mathcal O})^*,&\text{所有 }i,
\end{cases}
\]
其中行列式视为 $\widetilde{\mathcal O}/\pi^h\widetilde{\mathcal O}$ 的元素。用公式（2.1）定义 $F:G_h\to G_h$。$G_h$ 是 $\widetilde k$ 上的代数群，$F:G_h\to G_h$ 是 $G_h$ 的某个 $k$ 有理结构（$k$-rational structure）的 Frobenius 映射。令 $T_h$ 为 $G_h$ 的对角矩阵群，$U_h$（相应地，$U_h^-$）为 $G_h$ 中对角线上均为一的上（下）三角矩阵群。我们可将 $X_h$ 等同于簇
\[
\{A\in G_h\mid A^{-1}F(A)\in U_h\cap FU_h^-\}.
\]
$G_h^F$ 通过左乘作用于 $X_h$，$T_h^F$ 通过右乘作用于 $X_h$。（注意 $\widetilde G^F=\varprojlim_hG_h^F$，$\widetilde T^F=\varprojlim_hT_h^F$。）关键是以下引理。
\begin{originaltheorem}{引理（Lemma）}
当 $h\geq2$ 时，$K_h=\ker(T_h\to T_{h-1})$\footnote{\textbf{译者注（本次校订）：}这里应为有限群的核 $K_h=\ker(T_h^F\to T_{h-1}^F)$。前面定义的右乘作用属于 $T_h^F$：对一般 $t\in T_h$，$(At)^{-1}F(At)=t^{-1}(A^{-1}F(A))F(t)$ 会多出对角因子 $t^{-1}F(t)$，因而不保持 $X_h$。后面取 $K_h$ 不变同调及对纤维求商时，均应使用这个有限核。} 在 $X_h$ 上的作用保持自然映射 $X_h\to X_{h-1}$ 的每个纤维（fibre），且 $X_h/K_h\to X_{h-1}$ 的每个纤维同构于 $\widetilde k$ 上的 $(n-1)$ 维仿射空间（affine space）。
\end{originaltheorem}
对于 $\widetilde k$ 上纯维数为 $d$ 的任一光滑代数簇 $Y$，定义
\[
H_i(Y)=H_c^{2d-i}(Y,\overline\Q_\ell(d)),
\]
其中 $\ell$ 是固定的、不等于 $\operatorname{char}\widetilde k$ 的素数。前面的引理表明，当 $h\geq2$ 时，$H_i(X_{h-1})$ 典范同构于 $H_i(X_h)^{K_h}$（$K_h$ 在 $H_i(X_h)$ 上的不动点）。特别地，得到自然嵌入 $H_i(X_{h-1})\to H_i(X_h)$。利用这些嵌入取归纳极限 $\varinjlim_hH_i(X_h)$，并将 $H_i(X)$ 定义为此极限。$H_i(X)$ 自然分解为直和 $\bigoplus H_i(X)_\theta$，其中 $\theta$ 遍历 $(\widetilde T^F)^\vee$。对固定的 $\theta$，$H_i(X)_\theta$ 有限维，并在 $i$ 足够大时为零。同时，它自然构成 $\widetilde G^F$ 模（$\widetilde G^F$ 的作用通过一个有限商分解）。令 $R_\theta=\sum_i(-1)^iH_i(X)_\theta$。
\begin{originaltheorem}{定理（Theorem）}
对每个 $\theta\in(\widetilde T^F)^\vee$，$\pm R_\theta$ 是不可约 $\widetilde G^F$ 模。若 $\theta\ne\theta'$，则 $\pm R_\theta$ 与 $\pm R_{\theta'}$ 不同。
\end{originaltheorem}
为证明该定理，只需令 $\Sigma=G_h^F\backslash(X_h\times X_h)$（$G_h^F$ 在 $X_h\times X_h$ 上作对角作用），并证明
\[
\sum_i(-1)^i\dim H_i(\Sigma)_{\theta,\theta'^{-1}}
=\begin{cases}1,&\theta'=\theta,\\0,&\theta'\ne\theta.\end{cases}
\]
（这里 $T_h^F\times T_h^F$ 通过右乘作用于 $\Sigma$，$H_i(\Sigma)_{\theta,\theta'^{-1}}$ 表示 $T_h^F\times T_h^F$ 的 $(\theta,\theta'^{-1})$ 特征空间。）

为计算 $\Sigma$ 的（等变）Euler 示性数（equivariant Euler characteristic），我们使用如下思想：在适当条件下，可通过向量场的零点计算空间的 Euler 示性数。

映射 $(g,g')\mapsto(x,x',y)$，其中 $x=g^{-1}F(g)$、$x'=g'^{-1}F(g')$、$y=g^{-1}g'$，定义一个同构
\[
\Sigma\simeq\{(x,x',y)\in(U_h\cap FU_h^-)\times(U_h\cap FU_h^-)\times G_h\mid xF(y)=yx'\}
\]
（比较 [1, 6.6]）。现在，每个 $y\in G_h$ 都可唯一写成
\[
\begin{aligned}
y&=y'_1y'_2y''_1y''_2,\\
y'_1&\in U_h^-\cap F^{-1}U_h,\qquad y'_2\in T_h(U_h^-\cap F^{-1}U_h^-),\\
y''_1&\in U_h\cap F^{-1}U_h^-,\qquad y''_2\in U_h\cap F^{-1}U_h.
\end{aligned}
\]
作变量代换 $xF(y'_1)=\widetilde x\in U_h$，则 $\Sigma$ 的方程变为
\[
\widetilde xF(y'_2)F(y''_1)F(y''_2)=y'_1y'_2y''_1y''_2x'.
\]
每个 $z\in U_h$ 都可唯一写成 $z=y''_2x'F(y''_2)^{-1}$，其中 $y''_2\in U_h\cap F^{-1}U_h$、$x'\in U_h\cap FU_h^-$。故 $\Sigma$ 的方程成为
\[
\widetilde xF(y'_2)F(y''_1)\in y'_1y'_2y''_1U_h=y'_1y'_2U_h.
\]
因此可将 $\Sigma$ 写成：
\[
\begin{aligned}
\Sigma=\{&(\widetilde x,y'_1,y'_2,y''_1)\in(U_h\cap FU_h^-)\times(U_h^-\cap F^{-1}U_h)\\
&\times T_h(U_h^-\cap F^{-1}U_h^-)\times(U_h\cap F^{-1}U_h^-)\mid\\
&\hspace{2em}(y'_2)^{-1}(y'_1)^{-1}\widetilde xF(y'_2)F(y''_1)\in U_h\}.
\end{aligned}
\]
$T_h^F\times T_h^F$ 的作用由下式给出：
\[
(t,t'):(\widetilde x,y'_1,y'_2,y''_1)\longmapsto
(t\widetilde xt^{-1},ty'_1t^{-1},ty'_2t'^{-1},t'y''_1t'^{-1}).
\]
这一作用可以延拓为一个更大的群 $H$ 的作用。$H$ 由所有 $n\times n$ 对角矩阵对 $(t,t')$ 构成，对角元属于 $(\widetilde{\mathcal O}/\pi^h\widetilde{\mathcal O})^*$，满足 $\det t=\det t'$，以及
\[
t^{-1}F(t)=t'^{-1}F(t')\in Z_{\GL_n}(T_h(U_h^-\cap F^{-1}U_h^-)).
\]
$H$ 的作用为
\[
\begin{aligned}
(t,t'):(\widetilde x,y'_1,y'_2,y''_1)\longmapsto
(&F(t)\widetilde xF(t)^{-1},F(t)y'_1F(t)^{-1},\\
&ty'_2t'^{-1},t'y''_1t'^{-1}).
\end{aligned}
\]
形如 $\diag(1,\ldots,1,a)$，$a\in(\widetilde{\mathcal O}/\pi^h\widetilde{\mathcal O})^*$ 的对角矩阵均中心化 $T_h(U_h^-\cap F^{-1}U_h^-)$。因此，当
\[
t=\diag(\xi^{\phi^{n-1}},\ldots,\xi^\phi,\xi)
\]
时，任意一对 $(t,t)$ 都属于 $H$；这里 $\xi\in(\widetilde{\mathcal O}/\pi^h\widetilde{\mathcal O})^*$ 是阶与 $\operatorname{char}k$ 互素的单位根。所有这些 $(t,t)$ 构成 $H$ 的一个一维子环面（one dimensional subtorus）$\mathscr T$（在 $\widetilde k$ 上）。此环面作用于 $\Sigma$，并与 $T_h^F\times T_h^F$ 的作用交换。它在 $\Sigma$ 上的不动点集是由 $\widetilde x=y'_1=y''_1=e$、$y'_2\in T_h^F$ 给出的有限集；因而它同构于 $T_h^F$，$T_h^F\times T_h^F$ 通过左右乘作用于其上。因此
\[
\begin{aligned}
\sum_i(-1)^i\dim H_i(\Sigma)_{\theta,\theta'^{-1}}
&=\sum_i(-1)^i\dim H_i(\Sigma^{\mathscr T})_{\theta,\theta'^{-1}}\\
&=\dim H_0(T_h^F)_{\theta,\theta'^{-1}}\\
&=\begin{cases}1,&\theta=\theta',\\0,&\theta\ne\theta',\end{cases}
\end{aligned}
\]
所需等式得证。

\origsection{3}{$\SL_2$ 的情形}
设 $V$ 是 $K$ 上的二维向量空间，令 $\widetilde V=V\otimes_k\widetilde K$\footnote{\textbf{译者注（本次校订）：}张量积的基域应为 $K$，即 $\widetilde V=V\otimes_K\widetilde K$。这里 $V$ 是 $K$ 向量空间，$k$ 是剩余域；在混合特征情形尤其不存在所需的自然 $k$ 代数结构。}，以 $F=1\otimes\phi$ 定义 $F:\widetilde V\to\widetilde V$。假设给定非零元素 $\omega\in\bigwedge^2V$。集合
\[
X=\{x\in\widetilde V\mid x\wedge Fx=\omega\}
\]
在 $\SL_2(V)$ 的自然作用下不变，也在群 $\mathsf T=\{\lambda\in\widetilde K^*\mid\lambda\lambda^\phi=1\}$ 的标量乘法作用下不变。（这个集合 $X$ 可等同于第 1 节中对 $G=\SL_2$、$T$ 为与 $K$ 的非分歧二次扩张相联系的极大环面所定义的 $X$。）与这一 $X$ 类似的对象出现在 Drinfeld 的工作中。下面说明在这个例子中如何定义同调群 $H_i(X)$。容易检验，若 $x\in X$，由 $x,Fx$ 张成的 $\widetilde{\mathcal O}$ 模在 $F$ 下不变，因而可由 $V$ 中的一个格作标量扩张得到（lattice）$L$。

因此 $X$ 是一个不交并 $X=\coprod_LX_L$，其中 $L$ 遍历 $V$ 中行列式为 $\mathcal O^*\omega$ 的所有格，且
\[
X_L=\{x\in L\otimes_{\mathcal O}\widetilde{\mathcal O}\mid x\wedge Fx=\omega\}.
\]
每个 $X_L$ 自然可视为 $\varprojlim_hX_{L,h}$，其中 $X_{L,h}$ 是由下式定义的 $\widetilde k$ 上 $h$ 维簇：
\[
X_{L,h}=\{x\in L\otimes_{\mathcal O}(\widetilde{\mathcal O}/\pi^h\widetilde{\mathcal O})\mid x\wedge Fx=\omega\}.
\]
$X_{L,h}\to X_{L,h-1}$ 的纤维具有类似第 2 节引理中的性质。因此可以仿照第 2 节定义 $H_i(X_L)=\varinjlim_hH_i(X_{L,h})$。再定义 $H_i(X)=\bigoplus_LH_i(X_L)$。

类似的论证应当也适用于一般情形。

\origsection{4}{有限层上的虚表示}
设 $G,T,B,U$ 如第 1 节。假设 $G$ 可由 $\Z$ 上的 Chevalley 群作标量扩张（extension of scalars）得到，从而可以定义 $G(\mathcal O)$ 与 $G(\widetilde{\mathcal O})$。令 $\widetilde G_h=G(\widetilde{\mathcal O}/\pi^h\widetilde{\mathcal O})$。假设 $T(\widetilde K)\subset G(\widetilde{\mathcal O})$。\footnote{\textbf{译者注（本次校订）：}这里须使用环面的整数点（最大有界子群），而不能使用整个 $T(\widetilde K)$。在本节的分裂情形及第1节的 Borel 子群假设下，$T$ 在 $\widetilde K$ 上分裂，非平凡余特征会给出含任意赋值的元素，所以整个 $T(\widetilde K)$ 不可能包含在 $G(\widetilde{\mathcal O})$ 中。相应地，$\widetilde T_h$ 应取 $T(\widetilde K)\cap G(\widetilde{\mathcal O})$ 的约化像；若已选定相容的整数模型，可记这一子群为 $T(\widetilde{\mathcal O})$。}令 $\widetilde T_h,\widetilde U_h$ 分别是 $T(\widetilde K)$、$U(\widetilde K)\cap G(\widetilde{\mathcal O})$ 在 $G(\widetilde{\mathcal O})\to\widetilde G_h$ 下的像。$F:G(\widetilde K)\to G(\widetilde K)$ 诱导 $F:\widetilde G_h\to\widetilde G_h$。这一映射为 $\widetilde k$ 上的代数群 $\widetilde G_h$ 给出一个 $k$ 有理结构。定义
\[
X_h=\{g\in\widetilde G_h\mid g^{-1}F(g)\in\widetilde U_h\}/(\widetilde U_h\cap F^{-1}\widetilde U_h).
\]
有限群 $\widetilde G_h^F\times\widetilde T_h^F$ 像前面一样，通过左乘与右乘作用于 $X_h$。对每个特征 $\theta:\widetilde T_h^F\to\overline\Q_\ell^*$，构造
\[
R_\theta=\sum_i(-1)^iH_i(X_h)_\theta.
\]
\begin{originaltheorem}{定理（Theorem）}
若 $\theta$ 足够正则，则虚 $\widetilde G_h^F$ 模（virtual module）$\pm R_\theta$ 不可约。它与 $B$ 的选择无关。
\end{originaltheorem}
$h=1$ 的情形由 [1] 得出。
\begin{thebibliography}{9}
\bibitem[1]{c09-1} P. Deligne and G. Lusztig, \emph{Representations of reductive groups over finite fields}, Ann. of Math. (2) 103 (1976), 103--161.
\bibitem[2]{c09-2} P. Gérardin, \emph{Cuspidal unramified series for central simple algebras over local fields}, 本论文集，第 1 册，第 157--169 页。
\end{thebibliography}
\par\noindent 英国 Coventry，Warwick 大学（University of Warwick）。
